%%%%%%%%%%%%%%%%%%%%%%%%%%%%%%%%%%%%%%%%%%%%%%%%%%%%%%%%%%%%%%%%%%%%%%%%
%                                                                      %
%     File: Thesis_Conclusions.tex                                     %
%     Tex Master: Thesis.tex                                           %
%                                                                      %
%     Author: Andre C. Marta                                           %
%     Last modified :  24 Jul 2026                                     %
%                                                                      %
%%%%%%%%%%%%%%%%%%%%%%%%%%%%%%%%%%%%%%%%%%%%%%%%%%%%%%%%%%%%%%%%%%%%%%%%

\chapter{Conclusions}
\label{chapter:conclusions}

This chapter concludes the dissertation, summarizing the key achievements and proposing directions for future research.

\section{Conclusions}
\label{sec:concl_summary}
% TODO: Summarize the findings of this thesis:
% - Review the development of the cooperative UAV-USV model and NMPC formulation.
% - Reflect on the C++ and ROS 2 implementation and solver performance in real-time.
% - Review simulation achievements in Gazebo Harmonic with MoorDyn.
% - State the final conclusions regarding the safety, feasibility, and performance of tethered marine-aerial operations.

\section{Future Work}
\label{sec:concl_future}
% TODO: Propose improvements and future lines of research:
% - Experimental validation: deploying on the physical x500 quadrotor and WAM-V catamaran.
% - Hardware-In-the-Loop (HIL) testing to validate PX4 and ROS 2 controller interfaces under real flight constraints.
% - Advanced wind and wave estimation techniques to feed the NMPC parameters.
% - Decentralized control architectures to improve scalability and communication robustness.
% - Energy harvesting: exploring tether power delivery optimization and winch regeneration.
